\documentclass{article}
\usepackage[utf8]{inputenc}
\usepackage[margin=2cm,includefoot]{geometry}
\usepackage{booktabs}
\usepackage{graphicx}
\usepackage{amsmath, amssymb, amsthm}
\usepackage{physics}
\usepackage[shortlabels]{enumitem}
\usepackage{cancel}
\usepackage{lastpage}
\usepackage{braket}
\usepackage{mathrsfs}
\usepackage{bm}
\usepackage{subcaption}
\usepackage{float}

\usepackage{fancyhdr}
\pagestyle{fancy}
\newcommand{\AssignmentNum}{2}
\lhead{Stromingsleer en transportverschijnselen (NS-265B)\\
2025/2026}
\rhead{Tereza Bílková (7137257), Bela F. Brunner (8002193),\\ Niels Epema (2927578), Adam Zich (0074187)}
\fancyfoot{}
\fancyfoot[R]{Page \thepage /\pageref{LastPage}}

\newcommand{\vba}{\mathbf{A}}
\newcommand{\vbb}{\mathbf{B}}
\newcommand{\C}{\mathbf{C}}
\newcommand{\x}{\mathbf{\hat{x}}}
\newcommand{\y}{\mathbf{\hat{y}}}
\newcommand{\z}{\mathbf{\hat{z}}}
\newcommand{\vr}{\mathbf{\hat r}}
\newcommand{\thet}{\boldsymbol{{\hat \theta}}}
\newcommand{\ph}{\boldsymbol{{\hat \phi}}}
\newcommand{\ba}{\mathbf{a}}
\newcommand{\bb}{\mathbf{b}}
\newcommand{\vl}{\mathbf{l}}
\newcommand{\vv}{\mathbf{v}}

\newcommand{\R}{\mathbb{R}}
\newcommand{\p}{\mathbb{P}}
\newcommand{\N}{\mathbb{N}}
\newcommand{\Z}{\mathbb{Z}}
\newcommand{\Q}{\mathbb{Q}}
\newcommand{\I}{\mathbb{I}}
\newcommand{\E}{\mathbb{E}}

\newcommand{\mbeq}{\overset{!}{=}}
\newcommand{\PATH}{figures/}

\begin{document}
    \setlength{\headheight}{23.10004pt}
    \begin{center}
    \LARGE{\bfseries{ Project \AssignmentNum}}
    \line(1,0){430}
    \end{center}

\section{Analytic part}

\subsection*{a)}

We argue that we can describe the flow using mass and momentum balance as

\begin{align*}
    \partial_x^2(\mathbf v)+\partial _x\nabla v_x&=\nabla ^2v_x,\\
    \partial ^2_xv_x+\partial ^2_xv_y+\partial^2_xv_z+\partial _x\nabla \cdot v_x&=\partial_x^2v_x+\partial_y^2v_x+\partial_z ^2v_x,\\
    \partial ^2_xv_x+\partial ^2_xv_y+\partial^2_xv_z+\partial _{xy}v_x+\partial_{xz}v_x&=\partial_y^2v_x+\partial_z ^2v_x,\\
    \partial^2_x(v_x+v_y+v_z)&=(\partial _y^2+\partial_z^2-\partial _{xy}-\partial_{xz})v_x.
\end{align*}

This is because TODO

\subsection*{b)}

In our situation, the equations, using $v_x = u$, simplify to

\begin{align*}
    \rho \,\partial_t u &= -\partial_x p + \mu \,\partial_z^2 u, \\
    \partial_z v_z &= 0, \qquad \text{hence } v_z = 0, \\
    0 &= -\partial_y p, \\
    0 &= -\partial_z p - \rho g, \\
    p(x,z,t) &= -\rho g z + C(x,y,t).
\end{align*}

This is because TODO

Additionally, $C$ does not depend on $x, y$, therefore we can siplify the equations even further to TODO

\subsection*{c)}

The equations of motion and boundary conditions then become
\begin{align*}
    (1)\quad & \rho \,\partial_t u = \mu \,\partial_z^2 u, \\
    (2)\quad & z = 0 :\; u = U(t), \\
    (3)\quad & z = \infty :\; u \rightarrow 0, \\
    (4)\quad & t = 0 :\; u = 0.
\end{align*}

This is beacause TODO

\subsection*{d)}
We can find with BC (2) that the solution should have a form: 
\begin{align*}
    u(z,t) = f(z)\cdot U_0 \sin\!\left( 2\pi \frac{t - \tau(z)}{T} \right),
\end{align*}
with $f(0)=1$ and $\tau(0)=0$. Here, $f(z)$ is a function that describes how the amplitude of the oscillations changes with height and $\tau(z)$ is a function that describes how the phase of the oscillations changes with height. Since we know that the amplitude of oscilation is damped for higher $z$, we can expect that $f(z)$ is a decreasing function of $z$. Furtermore, Damping is usually exponential, which is why $f(z)=e^{-k}$ is a reasonable guess. The phase-shift is necessary, since we know that the oscilation only moves up the water with a ceratin delay. Then, the equaiton of motion becomes
\begin{align*}
    u(z,t) = U_0 e^{-kz} \sin\!\left( 2\pi \frac{t - \tau(z)}{T} \right),
\end{align*}
where $k$ must be a positive constant, since otherwise there would be no damping effect.\\

Note that we implicitly used BC (3), as this is determining the damping. Furthermore, this is only reasonable after a few oscilations, which is why we cannot apply BC (4). This leaves us with BC (1), with which we can determine $k$ and $\tau(z)$:
\begin{align*}
    \rho \,\partial_t u &= \mu \,\partial_z^2 u\\
    \rho U_0 e^{-kz} \frac{2\pi}{T} \cos\!\left( 2\pi \frac{t - \tau(z)}{T} \right) &=\mu\partial _z\left(U_0e^{-kz}(-k)\sin\left(2\pi\frac{t-\tau(z)}{T}\right)-U_0e^{-kz}\cos\left(2\pi\frac{t-\tau(z)}{T}\right)\frac{2\pi\partial_z\tau(z)}{T}\right)\\
    \frac{\rho}{\mu}e^{-kz}\frac{2\pi}{T}\cos\!\left( 2\pi \frac{t - \tau(z)}{T} \right) &=e^{-kz}\left(k^2\sin \left(2\pi\frac{t-\tau(z)}{T}\right)+k\cos \left(2\pi\frac{t-\tau(z)}{T}\right)\frac{2\pi\partial_z\tau(z)}{T}\right)\\
    &\quad -e^{-kz}\cos\left(2\pi\frac{t-\tau(z)}{T}\right)\frac{2\pi}{T}\left((-k\partial _z\tau(z))+\partial ^2_z\tau(z)\right)\\
    &\quad-e^{-kz}\sin\left(2\pi\frac{t-\tau(z)}{T}\right)\left(\frac{2\pi\partial_z\tau(z)}{T}\right)^2\\
    \frac{2\pi\rho}{\mu T}\cos\!\left( 2\pi \frac{t - \tau(z)}{T} \right) &=\sin \left(2\pi\frac{t-\tau(z)}{T}\right)\left(k^2-\left(\frac{2\pi\partial_z\tau(z)}{T}\right)^2\right)\\
    &\quad+\cos \left(2\pi\frac{t-\tau(z)}{T}\right)\left(2k\frac{2\pi\partial_z\tau(z)}{T}-\left(\frac{2\pi\partial_z^2\tau(z)}{T}\right)\right)
\end{align*}
Since we know that this must be true for all of $t$, we can find that the $\sin$ terms must be zero and the $\cos$ terms must be equal on both sides. This gives us the two equations: 
\begin{align*}
    k^2-\left(\frac{2\pi\partial_z\tau(z)}{T}\right)^2&=0,\\
    2k\frac{2\pi\partial_z\tau(z)}{T}-\left(\frac{2\pi\partial_z^2\tau(z)}{T}\right)&=\frac{2\pi\rho}{\mu T}.
\end{align*}
The first one reduces to: 
\begin{align*}
    k&=\frac{2\pi\partial_z\tau(z)}{T} \Leftrightarrow \partial_z\tau(z) = \frac{kT}{2\pi} \rightarrow \tau(z) = \frac{kT}{2\pi}z + C
\end{align*}
To determine the constant, we use BC (2):\footnote{Of course, used this already above, for completion the calculation is included here.}
\begin{align*}
    u(0,t)&=U(t)\\
    U_0 \sin\!\left( 2\pi \frac{t-\tau(0)}{T} \right)&=  U_0 \sin\!\left( 2\pi \frac{t}{T} \right)\\
    \Rightarrow \tau (0)&=0 \Rightarrow \frac{kT}{2\pi}0 + C = 0 \Rightarrow C = 0
\end{align*}
With the second equation, we can determine the value of $k$:
\begin{align*}
    2k\partial_z\tau(z)-\partial_z^2\tau(z)&=\frac{\rho}{\mu}\\
    \Leftrightarrow\tau''-2k\tau'+\frac{\rho}{\mu}&=0\\
    0-2k\frac{kT}{2\pi}+\frac{\rho}{\mu}&=0\\
    \Rightarrow k&=\sqrt{\frac{\rho 2\pi}{2\mu T}}=\sqrt{\frac{\rho \omega}{2\mu}}
\end{align*}

We find that $k = \sqrt{\dfrac{\rho \omega}{2\mu}}$ where $\omega = 2\pi/T$ and $\tau(z) = \frac{kT}{2\pi}z$. \\
The statement that $k$ increases with $\omega$ is equivalent to saying that a faster oscilation will experience stronger damping. This statement is reasonable, since the acceleration of the next layer of fluid due to the friction is not instentanous. However, when now the oscilation is faster, the next layer of fluid will not have enough time to accelerate before the oscilation changes direction again, therefore the whole motion experiences stronger damping per height.

\subsection*{e)}

The velocity field $u(z, t)$ changes over a period $T$ as a damped sinusodial function along the z-axis. \\

When we plot the field at time $t = 0$ (Figure \ref{fig:velocity_t0}) and $t = T/4$ (Figure \ref{fig:velocity_tT/4}) we see that TODO

    \begin{figure}[H]
    \centering
    
    \begin{subfigure}{0.48\textwidth}
        \centering
        \includegraphics[width=\linewidth]{\PATH 1_e_Velocity_profile_t=0.png}
        \caption{}
        \label{fig:velocity_t0}
    \end{subfigure}
    \hfill
    \begin{subfigure}{0.48\textwidth}
        \centering
        \includegraphics[width=\linewidth]{\PATH 1_e_Velocity_profile_t=Tover4.png}
        \caption{}
        \label{fig:velocity_tT/4}
    \end{subfigure}

    \caption{Velocity field for $t = 0$ (a) and $t = T/4$ (b)}
    \end{figure}

\section{Computational part}

All explanations for the computational part are in the jupyter notebook.

\end{document}